\begin{trapbox}

	\begin{description}[style=nextline, leftmargin=2.8em, labelsep=0.5em, itemsep=0.35em]
		\item[\textbf{\textcolor{CTrap}{易错点一（忽视平行条件）}}]
			直接取直线上一点到平面的距离作为线面距，但未确认直线与平面平行；若不平行，距离不是定值。
		\item[\textbf{\textcolor{CTrap}{正确理解（先判平行）：}}]
			只有当直线与平面平行时，线面距才等于点到平面的距离，且该距离与点选择无关。
		\item[\textbf{\textcolor{CTrap}{错因分析（条件遗漏）：}}]
			忽略结论的前提，将适用范围误认为所有情况。
	\end{description}

	\medskip
	\begin{description}[style=nextline, leftmargin=2.8em, labelsep=0.5em, itemsep=0.35em]
		\item[\textbf{\textcolor{CTrap}{易错点二（公式漏项）}}]
			计算距离时只写分子 $|Ax_0+By_0+Cz_0+D|$，忘记除以 $\sqrt{A^2+B^2+C^2}$。
		\item[\textbf{\textcolor{CTrap}{正确理解（公式完整）：}}]
			距离公式包含分母 $\sqrt{A^2+B^2+C^2}$，且分子要取绝对值。
		\item[\textbf{\textcolor{CTrap}{错因分析（记忆偏差）：}}]
			点面距公式为绝对值除以模长，部分学生只记住分子形式。
	\end{description}

	\medskip
	\begin{description}[style=nextline, leftmargin=2.8em, labelsep=0.5em, itemsep=0.35em]
		\item[\textbf{\textcolor{CTrap}{易错点三（代错平面）}}]
			求两平行平面 $\alpha$ 和 $\beta$ 的距离时，在 $\alpha$ 上取点代入 $\alpha$ 的方程，得到 0，误以为距离为 0。
		\item[\textbf{\textcolor{CTrap}{正确理解（代对平面）：}}]
			应取一个平面上的点，代入另一个平面的方程计算距离。
		\item[\textbf{\textcolor{CTrap}{错因分析（公式理解不清）：}}]
			点面距公式中的平面是目标平面，点应来自另一个平面或几何体。
	\end{description}

\end{trapbox}
